\documentclass[10pt,a4paper]{amsart}
\usepackage[margin=19mm]{geometry}
\usepackage{amsmath,amssymb,mathtools}
\usepackage[scr=rsfs,cal=euler]{mathalfa}
\usepackage{xfrac}
\usepackage[unicode,hidelinks,bookmarks=false]{hyperref}
\newcommand{\ie}{i.e.\ }
\newcommand{\cf}{cf.\ }
\newcommand{\resp}{respectively\ }
\newcommand{\Subsection}{\subsection}
\providecommand{\nobreakdash}{\nobreak\hbox{-}}
\setlength{\emergencystretch}{2em}
\multlinegap0pt
\usepackage{bm}
\usepackage{bbm}
\usepackage{dsfont}
\usepackage{xspace}
\usepackage{cancel}
\usepackage{mathtools}
\usepackage{amsthm}
\makeatletter
\@ifundefined{lemma}{\newtheorem{lemma}{Lemma}}{}
\makeatother
\newcommand{\floor}[1]{\left\lfloor #1\right\rfloor}
\title[Adjacent comparison bounds and extremal sets for Ruzsa numbe]{Adjacent comparison bounds and extremal sets for Ruzsa numbers}
\author{Yuchen Ding, Huixi Li, Junfeng Li, Wei Niu, Xiamiao Zhao}
\date{Source: arXiv:2606.11069v1 (CC BY 4.0)}
\begin{document}
\maketitle

\noindent\emph{Source and licence.} Adjacent comparison bounds and extremal sets for Ruzsa numbers. Version arXiv:2606.11069v1 (CC BY 4.0), distributed under the Creative Commons Attribution 4.0 International licence, as recorded in the arXiv licence metadata for this exact version and shown on the abstract page https://arxiv.org/abs/2606.11069v1. Full rights notice: licenses/arxiv-2606.11069-CC-BY-4.0.txt.\par\smallskip

\noindent\textit{Editorial scope.} Counted unit: the Lemma whose source label is \texttt{lem:1} in the article (source0.tex lines 627--630, source label \texttt{lem:1}), delivered together with the article's own complete original proof of it (source0.tex lines 631--864, one \texttt{proof} environment of 234 lines). No mathematics is written, repaired or re-derived: statement and proof are the article's own text line for line. Required context retained uncounted: none. Every definition, display and label of those lines is retained unchanged. Omitted from this package and not redistributed: 1--626, 865--1330. Nothing inside a retained excerpt is omitted or altered: every retained body is delivered exactly as the article's lines are written, so no omission receipt of any kind accompanies this package. Citations: the retained excerpts carry no citation command at all, so no bibliography entry of the article is transcribed and no reference is redistributed. Reference adaptation: none was needed and none was made; every reference inside the delivered unit resolves inside it, so no unresolved reference or citation is produced. Printed slips kept literal: no printed word, symbol, hypothesis or number of the counted statement or of its proof was altered. Layout and class: the article's own class is \texttt{amsart} with packages including lineno, amsmath, amssymb, amsfonts, amsthm, latexsym, enumerate, url, cases, mathrsfs, array, longtable, xcolor, hyperref; the wrapper is the lane's pinned \texttt{amsart} editorial wrapper at 10pt on A4, which restates the theorem-like environments the retained text needs and adds the article's own macro definitions that the retained text needs (\texttt{\textbackslash floor}), with extra wrapper packages (bm, bbm, dsfont, xspace, cancel, mathtools). The delivered numbering is the unit's own. Assets: no retained span contains a figure environment, a graphics-inclusion command, a PDF-page inclusion, a table or a TikZ picture; no image file is copied into this package, the package declares no asset dependency and no image file is redistributed with it. Licence: version arXiv:2606.11069v1 of this article is distributed under the Creative Commons Attribution 4.0 International licence, as recorded in the arXiv licence metadata for this exact version and shown on the abstract page https://arxiv.org/abs/2606.11069v1. A full rights notice is delivered at licenses/arxiv-2606.11069-CC-BY-4.0.txt. Characters: every retained excerpt is pure ASCII, so the delivered engine, XeLaTeX with its default Latin Modern text font, reports no missing character.
 \begin{lemma}\label{lem:1}
   For any positive integer $m$ we have both
   $$R_m\le 4R_{m+1} \quad \text{and} \quad R_{m+1}\le 4R_{m}.$$
 \end{lemma}

 \begin{proof}
 {\bf We first prove the assertion $R_{m+1}\le 4R_{m}$}. 

 Suppose that $\mathcal{A}$ is a subset of $\mathbb{Z}_m$ such that for any $n\in \mathbb{Z}_m$ we have
 $$
 1\le \sigma_{\mathcal{A}}(n)\le R_{m}.
 $$
Without loss of generality, we may assume $0\le a\le m-1$ for any $a\in \mathcal{A}$. Define
$$
\mathcal{A}_{1}=\big\{a\in \mathcal{A}: a\ge \floor{m/2}\big\}
$$
and
$$
\mathcal{A}_{+}=\mathcal{A}_{1}+1:=\{a+1: a\in \mathcal{A}_{1}\}.
$$
Now, we let
$
\mathcal{B}=\mathcal{A}\cup \mathcal{A}_{+},
$
and regard $\mathcal{B}$ as a subset of $\mathbb{Z}_{m+1}$.

First, we show that $\mathcal{B}$ is an additive basis of $\mathbb{Z}_{m+1}$, i.e., $\sigma_{\mathcal{B}}(n)\ge 1$ for any $n\in \mathbb{Z}_{m+1}$. For any $0\le n\le m-1$, there are two elements $a, a'\in \mathcal{A}$ such that
$$
n\equiv a+a'\pmod{m},
$$
which implies $n=a+a'$ or $n=a+a'-m$ since $0\le a+a'\le 2m-2$. If $n=a+a'$, then 
$
n\equiv a+a'\pmod{m+1}.
$
Now, suppose that $n=a+a'-m$. Then $a+a'\ge m$, and hence at least one of $a$ and $a'$ is $\ge \floor{m/2}$. We may assume without loss of generality that $a\ge \floor{m/2}$. Then we have
$$
n=(a+1)+a'-(m+1)\equiv (a+1)+a'\pmod{m+1},
$$
which means that $\sigma_{\mathcal{B}}(n)\ge 1$ for any $0\le n\le m-1$. It remains to show $\sigma_{\mathcal{B}}(m)\ge 1$. Note that
there are two elements $a_1$ and $a_2$ of $\mathcal{A}$ such that
$$
m-1\equiv a_1+a_2\pmod{m}.
$$
Since $0\le a_1+a_2\le 2m-2$, we must have $m-1=a_1+a_2$. It follows that at least one of $a_1$ and $a_2$ is $\ge  \floor{m/2}$. Without loss of generality, we assume that $a_1\ge  \floor{m/2}$. Hence,
$$
m=(a_1+1)+a_2 \equiv (a_1+1)+a_2\pmod{m+1},
$$
completing the proof of $\sigma_{\mathcal{B}}(n)\ge 1$ for any $n\in \mathbb{Z}_{m+1}$.

Next, we will prove $\sigma_{\mathcal{B}}(n)\le 4R_m$ for any $0\le n\le m$. Let
$$
\mathcal{A}_{2}=\mathcal{A}\setminus\mathcal{A}_1=\big\{a\in \mathcal{A}: a< \floor{m/2}\big\}.
$$
By this symbol, we have
$$
\mathcal{B}=\mathcal{A}_1\cup \mathcal{A}_{2}\cup\mathcal{A}_{+}.
$$
For any subsets $\mathcal{S}_1$ and $\mathcal{S}_2$ of integers, define
$$
f_{\mathcal{S}_1,\mathcal{S}_2}(n)=\#\big\{(s_1,s_2)\in \mathcal{S}_1\times\mathcal{S}_2: n=s_1+s_2\big\}.
$$
We first consider the range $0\le n\le m-1$. Suppose that
\begin{align*}
    n\equiv b_1+b_2\pmod{m+1},
\end{align*}
where $b_1$ and $b_2$ are two elements of $\mathcal{B}$.
Then $n=b_1+b_2$ or $n=b_1+b_2-(m+1)$ since $0\le b_1+b_2\le 2m$. In the case $n=b_1+b_2$ we have
\begin{align*}
    \#\big\{(b_1, b_2)&\in \mathcal{B}^2 :  n=b_1 +b_2\big\}
    \le f_{\mathcal{A}_{1},\mathcal{A}_{1}}(n)+f_{\mathcal{A}_{1},\mathcal{A}_{2}}(n)+f_{\mathcal{A}_{2},\mathcal{A}_{1}}(n)\\
&+f_{\mathcal{A}_{1},\mathcal{A}_{+}}(n)+f_{\mathcal{A}_{+},\mathcal{A}_{1}}(n)+f_{\mathcal{A}_{2},\mathcal{A}_{2}}(n)+f_{\mathcal{A}_{2},\mathcal{A}_{+}}(n)+f_{\mathcal{A}_{+},\mathcal{A}_{2}}(n)+f_{\mathcal{A}_{+},\mathcal{A}_{+}}(n).
\end{align*}
Note that
$$
f_{\mathcal{A}_{1},\mathcal{A}_{+}}(n)=f_{\mathcal{A}_{1},\mathcal{A}_{1}}(n-1)=0,\quad f_{\mathcal{A}_{2},\mathcal{A}_{+}}(n)=f_{\mathcal{A}_{2},\mathcal{A}_{1}}(n-1),
$$
$$
f_{\mathcal{A}_{+},\mathcal{A}_{1}}(n)=f_{\mathcal{A}_{1},\mathcal{A}_{1}}(n-1)=0,\quad f_{\mathcal{A}_{+},\mathcal{A}_{2}}(n)=f_{\mathcal{A}_{1},\mathcal{A}_{2}}(n-1),
$$
and
$$
f_{\mathcal{A}_{+},\mathcal{A}_{+}}(n)=f_{\mathcal{A}_{1},\mathcal{A}_{1}}(n-2)=0
$$
since $a+a'\ge 2\floor{m/2}\ge m-1>n-1$ for any $a, a'\in \mathcal{A}_{1}$,
from which we obtain
\begin{align}\label{eq-proof-1}
    \#\big\{(b_1, b_2)\in \mathcal{B}^2& :  n=b_1 +b_2\big\}
    \le f_{\mathcal{A}_{1},\mathcal{A}_{1}}(n)+f_{\mathcal{A}_{1},\mathcal{A}_{2}}(n)\nonumber\\
    &+f_{\mathcal{A}_{2},\mathcal{A}_{1}}(n)+f_{\mathcal{A}_{2},\mathcal{A}_{2}}(n)+f_{\mathcal{A}_{2},\mathcal{A}_{1}}
    (n-1)+f_{\mathcal{A}_{1},\mathcal{A}_{2}}
    (n-1).
\end{align}
Similarly, in the case $n=b_1+b_2-(m+1)$ we have
\begin{align}\label{eq-proof-2}
    \#\big\{(b_1,b_2&)\in \mathcal{B}^2: n + (m+1)= b_1+b_2\big\}\nonumber\\
    &\le f_{\mathcal{A}_{1},\mathcal{A}_{1}}(n+1+m)+f_{\mathcal{A}_{1},\mathcal{A}_{2}}(n+1+m)
    +f_{\mathcal{A}_{2},\mathcal{A}_{1}}(n+1+m)\nonumber\\
    &\quad \quad +f_{\mathcal{A}_{1},\mathcal{A}_{+}}(n+1+m)+f_{\mathcal{A}_{+},\mathcal{A}_{1}}(n+1+m)+f_{\mathcal{A}_{2},\mathcal{A}_{2}}(n+1+m)\nonumber\\
    &\quad \quad +f_{\mathcal{A}_{2},\mathcal{A}_{+}}(n+1+m)+f_{\mathcal{A}_{+},\mathcal{A}_{2}}(n+1+m)+f_{\mathcal{A}_{+},\mathcal{A}_{+}}(n+1+m)\nonumber\\
    &= f_{\mathcal{A}_{1},\mathcal{A}_{1}}(n+1+m)+f_{\mathcal{A}_{1},\mathcal{A}_{2}}(n+1+m)
    +f_{\mathcal{A}_{2},\mathcal{A}_{1}}(n+m+1)\nonumber\\
    &\quad \quad +2f_{\mathcal{A}_{1},\mathcal{A}_{1}}(n+m)+f_{\mathcal{A}_{2},\mathcal{A}_{2}}(n+1+m)\nonumber\\
    &\quad \quad +f_{\mathcal{A}_{2},\mathcal{A}_{1}}(n+m)+f_{\mathcal{A}_{1},\mathcal{A}_{2}}(n+m)+f_{\mathcal{A}_{1},\mathcal{A}_{1}}(n-1+m).
\end{align}
Hence, we conclude from (\ref{eq-proof-1}) and (\ref{eq-proof-2}),
\begin{align}\label{eq-proof-3}
    \sigma_{\mathcal{B}}(n)&\le f_{\mathcal{A},\mathcal{A}}(n)+2f_{\mathcal{A},\mathcal{A}}(n+m)+f_{\mathcal{A},\mathcal{A}}(n-1)+f_{\mathcal{A},\mathcal{A}}(n-1+m)+f_{\mathcal{A},\mathcal{A}}(n+1+m)\nonumber\\
    &\le 2\#\big\{(a_1,a_2)\in \mathcal{A}^2: n\equiv a_1+a_2 \pmod{m}\big\}\nonumber\\
    &\quad \quad +\#\big\{(a_1,a_2)\in \mathcal{A}^2: n-1\equiv a_1+a_2\pmod{m}\big\}\nonumber\\
    &\quad \quad +\#\big\{(a_1,a_2)\in \mathcal{A}^2: n+1\equiv a_1+a_2\pmod{m}\big\}\nonumber\\
    &\le 4R_m.
\end{align}
It remains to estimate $\sigma_{\mathcal B}(m)$. Suppose that
$$
m\equiv b_1+b_2\pmod{m+1},
$$
where $b_1,b_2\in\mathcal B$. Since $0\le b_1+b_2\le 2m$, we must have $b_1+b_2=m$. Each $b_i$ can be written in the form $b_i=a_i+\varepsilon_i$, where $a_i\in\mathcal A$ and $\varepsilon_i\in\{0,1\}$; here $\varepsilon_i=1$ can occur only when $a_i\in\mathcal A_1$. Counting all such choices gives an upper bound, and hence
\begin{align*}
\sigma_{\mathcal B}(m)
&\le f_{\mathcal A,\mathcal A}(m)+2f_{\mathcal A,\mathcal A}(m-1)+f_{\mathcal A,\mathcal A}(m-2)\\
&\le \sigma_{\mathcal A}(0)+2\sigma_{\mathcal A}(-1)+\sigma_{\mathcal A}(-2)\\
&\le 4R_m.
\end{align*}
Therefore, we have $R_{m+1}\le 4 R_m$.

\medskip

{\bf Proof of the assertion $R_{m}\le 4R_{m+1}$}. 

The cases $m=1$ and $2$ are trivial. Now, we assume $m\ge 3$.
Suppose that $\mathcal{A}$ is a subset of $\mathbb{Z}_{m+1}$ such that for any $n\in \mathbb{Z}_{m+1}$ we have
 $$
 1\le \sigma_{\mathcal{A}}(n)\le R_{m+1}.
 $$
Without loss of generality, we may assume $0\le a\le m$ for any $a\in \mathcal{A}$. Define
$$
\mathcal{A}_{1}=\big\{a\in \mathcal{A}: a\ge m/2+1\big\}
$$
and
$$
\mathcal{A}_{-}=\mathcal{A}_{1}-1:=\{a-1: a\in \mathcal{A}_{1}\}.
$$
Now, we let
$
\mathcal{B}=\mathcal{A}\cup \mathcal{A}_{-},
$
and regard $\mathcal{B}$ as a subset of $\mathbb{Z}_{m}$.

First, we show that $\mathcal{B}$ is an additive basis of $\mathbb{Z}_{m}$, i.e., $\sigma_{\mathcal{B}}(n)\ge 1$ for any $n\in \mathbb{Z}_{m}$. We first consider $n=0$. Since $\mathcal A$ is an additive basis of $\mathbb Z_{m+1}$, there exist $a,a'\in\mathcal A$ such that
$$
m\equiv a+a'\pmod{m+1}.
$$
Since $0\le a+a'\le 2m$, we must have $a+a'=m$. Hence $0\equiv a+a'\pmod m$, and so $\sigma_{\mathcal B}(0)\ge1$.

Now let $1\le n\le m-1$. There are two elements $a, a'\in \mathcal{A}$ such that
$$
n\equiv a+a'\pmod{m+1},
$$
which implies $n=a+a'$ or $n=a+a'-(m+1)$ since $0\le a+a'\le 2m$. If $n=a+a'$, then 
$
n\equiv a+a'\pmod{m}.
$
Now, suppose that $n=a+a'-(m+1)$. Then $a+a'=n+m+1\ge m+2$. If both $a$ and $a'$ were $<m/2+1$, then, since $a$ and $a'$ are integers, we would have $a+a'\le m+1$, a contradiction. Thus at least one of $a$ and $a'$ is $\ge m/2+1$. We may assume without loss of generality that $a\ge m/2+1$. Then we have
$$
n=(a-1)+a'-m\equiv (a-1)+a'\pmod{m},
$$
which means that $\sigma_{\mathcal{B}}(n)\ge 1$ for any $1\le n\le m-1$. Hence $\mathcal B$ is an additive basis of $\mathbb Z_m$.

Next, we will prove $\sigma_{\mathcal{B}}(n)\le 4R_{m+1}$ for any $0\le n\le m-1$. Let
$$
\mathcal{A}_{2}=\mathcal{A}\setminus\mathcal{A}_1=\big\{a\in \mathcal{A}: a< m/2+1\big\}.
$$
By this symbol, we have
$$
\mathcal{B}=\mathcal{A}_1\cup \mathcal{A}_{2}\cup\mathcal{A}_{-}.
$$
It would be convenient to consider $\sigma_{\mathcal{B}}(0)$ separately. Suppose that
$$
b_1+b_2\equiv 0\pmod{m},
$$
where $b_1,b_2\in\mathcal B$. Then $b_1+b_2=0$ or $m$ or $2m$. The cases $b_1+b_2=0$ and $b_1+b_2=2m$ contribute at most two pairs in total. If $b_1+b_2=m$, then the corresponding elements of $\mathcal A$ have sum $m$, $m+1$, or $m+2$, according as neither, exactly one, or both of the two summands come from $\mathcal A_-$. Therefore
$$
\sigma_{\mathcal{B}}(0)\le \sigma_{\mathcal{A}}(-1)+\sigma_{\mathcal{A}}(0)+\sigma_{\mathcal{A}}(1)+2\le 3R_{m+1}+2\le 4R_{m+1},
$$
where we used $R_{m+1}\ge 2$.

In what follows, let $n$ be an integer satisfying $1\le n\le m-1$. Suppose that
\begin{align*}
    n\equiv b_1+b_2\pmod{m},
\end{align*}
where $b_1$ and $b_2$ are two elements of $\mathcal{B}$.
Then $n=b_1+b_2$ or $n=b_1+b_2-m$ since $0\le b_1+b_2\le 2m$. In the case $n=b_1+b_2$ we have
\begin{align*}
    \#\big\{(b_1, b_2)&\in \mathcal{B}^2 :  n=b_1 +b_2\big\}
    \le f_{\mathcal{A}_{1},\mathcal{A}_{1}}(n)+f_{\mathcal{A}_{1},\mathcal{A}_{2}}(n)+f_{\mathcal{A}_{2},\mathcal{A}_{1}}(n)\\
&+f_{\mathcal{A}_{1},\mathcal{A}_{-}}(n)+f_{\mathcal{A}_{-},\mathcal{A}_{1}}(n)+f_{\mathcal{A}_{2},\mathcal{A}_{2}}(n)+f_{\mathcal{A}_{2},\mathcal{A}_{-}}(n)+f_{\mathcal{A}_{-},\mathcal{A}_{2}}(n)+f_{\mathcal{A}_{-},\mathcal{A}_{-}}(n).
\end{align*}
Note that
$$
f_{\mathcal{A}_{1},\mathcal{A}_{-}}(n)=f_{\mathcal{A}_{1},\mathcal{A}_{1}}(n+1)=0,\quad f_{\mathcal{A}_{2},\mathcal{A}_{-}}(n)=f_{\mathcal{A}_{2},\mathcal{A}_{1}}(n+1),
$$
$$
f_{\mathcal{A}_{-},\mathcal{A}_{1}}(n)=f_{\mathcal{A}_{1},\mathcal{A}_{1}}(n+1)=0,\quad f_{\mathcal{A}_{-},\mathcal{A}_{2}}(n)=f_{\mathcal{A}_{1},\mathcal{A}_{2}}(n+1),
$$
and
$$
f_{\mathcal{A}_{-},\mathcal{A}_{-}}(n)=f_{\mathcal{A}_{1},\mathcal{A}_{1}}(n+2)=0
$$
since $a+a'\ge m+2>n+2$ for any $a, a'\in \mathcal{A}_{1}$,
from which we obtain
\begin{align}\label{eq-proof-4}
    \#\big\{(b_1, b_2)\in \mathcal{B}^2& :  n=b_1 +b_2\big\}
    \le f_{\mathcal{A}_{1},\mathcal{A}_{1}}(n)+f_{\mathcal{A}_{1},\mathcal{A}_{2}}(n)\nonumber\\
    &+f_{\mathcal{A}_{2},\mathcal{A}_{1}}(n)+f_{\mathcal{A}_{2},\mathcal{A}_{2}}(n)+f_{\mathcal{A}_{2},\mathcal{A}_{1}}
    (n+1)+f_{\mathcal{A}_{1},\mathcal{A}_{2}}
    (n+1).
\end{align}
Similarly, in the case $n=b_1+b_2-m$ we have
\begin{align}\label{eq-proof-5}
    \#\big\{(b_1,b_2&)\in \mathcal{B}^2: n +m= b_1+b_2\big\}\nonumber\\
    &\le f_{\mathcal{A}_{1},\mathcal{A}_{1}}(n+m)+f_{\mathcal{A}_{1},\mathcal{A}_{2}}(n+m)
    +f_{\mathcal{A}_{2},\mathcal{A}_{1}}(n+m)\nonumber\\
    &\quad \quad +f_{\mathcal{A}_{1},\mathcal{A}_{-}}(n+m)+f_{\mathcal{A}_{-},\mathcal{A}_{1}}(n+m)+f_{\mathcal{A}_{2},\mathcal{A}_{2}}(n+m)\nonumber\\
    &\quad \quad +f_{\mathcal{A}_{2},\mathcal{A}_{-}}(n+m)+f_{\mathcal{A}_{-},\mathcal{A}_{2}}(n+m)+f_{\mathcal{A}_{-},\mathcal{A}_{-}}(n+m)\nonumber\\
    &= f_{\mathcal{A}_{1},\mathcal{A}_{1}}(n+m)+f_{\mathcal{A}_{1},\mathcal{A}_{2}}(n+m)
    +f_{\mathcal{A}_{2},\mathcal{A}_{1}}(n+m)\nonumber\\
    &\quad \quad +2f_{\mathcal{A}_{1},\mathcal{A}_{1}}(n+m+1)+f_{\mathcal{A}_{2},\mathcal{A}_{2}}(n+m)\nonumber\\
    &\quad \quad +f_{\mathcal{A}_{2},\mathcal{A}_{1}}(n+m+1)+f_{\mathcal{A}_{1},\mathcal{A}_{2}}(n+m+1)+f_{\mathcal{A}_{1},\mathcal{A}_{1}}(n+2+m).
\end{align}
Hence, we conclude from (\ref{eq-proof-4}) and (\ref{eq-proof-5}),
\begin{align*}
    \sigma_{\mathcal{B}}(n)&\le f_{\mathcal{A},\mathcal{A}}(n)+2f_{\mathcal{A},\mathcal{A}}(n+1+m)+f_{\mathcal{A},\mathcal{A}}(n+1)+f_{\mathcal{A},\mathcal{A}}(n+2+m)+f_{\mathcal{A},\mathcal{A}}(n+m)\nonumber\\
    &\le2\#\big\{(a_1,a_2)\in \mathcal{A}^2: n\equiv a_1+a_2 \pmod{m+1}\big\}\nonumber\\
    &\quad \quad +\#\big\{(a_1,a_2)\in \mathcal{A}^2: n+1\equiv a_1+a_2\pmod{m+1}\big\}\nonumber\\
    &\quad \quad +\#\big\{(a_1,a_2)\in \mathcal{A}^2: n-1\equiv a_1+a_2\pmod{m+1}\big\}\nonumber\\
    &\le 4R_{m+1}.
\end{align*}
Therefore, we have $R_{m}\le 4 R_{m+1}$, completing the proof of our lemma.
 \end{proof}

\end{document}
